\documentclass{article}
\usepackage{amsmath,amssymb,amsthm}
\usepackage{algorithm}
\usepackage{algpseudocode}
\usepackage{enumitem}
\usepackage{hyperref}
\newtheorem{theorem}{Theorem}[section]
\newtheorem{lemma}{Lemma}[section]
\newtheorem{assumption}{Assumption}[section]

\begin{document}

%%%%%%%%%%%%%%%%%%%%%%%%%%%%%%%%%%%%%%%%%%%%%%%%%%%%%%%%%%%%
\section*{Algorithm Name}
\textbf{SVRG under the Polyak‑\L{}ojasiewicz (PL) condition}
%%%%%%%%%%%%%%%%%%%%%%%%%%%%%%%%%%%%%%%%%%%%%%%%%%%%%%%%%%%%

\section*{Mathematical Formulation}
Let 
\[
f(x)=\frac{1}{n}\sum_{i=1}^{n}f_i(x),\qquad 
f_i:\mathbb{R}^d\rightarrow\mathbb{R}
\]
be a finite‑sum objective.  
The algorithm proceeds in epochs $s=0,1,2,\dots$.
\begin{itemize}
    \item \textbf{Snapshot point.} At the beginning of epoch $s$ set
    \[
    \tilde{x}^{s}=x_{0}^{s}\in\mathbb{R}^{d}.
    \]
    \item \textbf{Full gradient at the snapshot.}
    \[
    \tilde{\mu}^{s}= \nabla f(\tilde{x}^{s})=\frac{1}{n}\sum_{i=1}^{n}\nabla f_i(\tilde{x}^{s}).
    \]
    \item \textbf{Inner loop.} For $t=0,\dots ,m-1$ draw an index $i_t$ uniformly
    from $\{1,\dots ,n\}$ and compute the \emph{variance‑reduced estimator}
    \[
    v_t^{s}= \nabla f_{i_t}\!\bigl(x_t^{s}\bigr)
            -\nabla f_{i_t}\!\bigl(\tilde{x}^{s}\bigr)+\tilde{\mu}^{s}.
    \]
    The iterate is updated by a standard SGD step
    \[
    x_{t+1}^{s}=x_{t}^{s}-\eta\,v_t^{s},
    \]
    where $\eta>0$ is a constant stepsize.
    \item \textbf{Epoch output.} After $m$ inner updates set
    \[
    \tilde{x}^{s+1}=x_{m}^{s}.
    \]
\end{itemize}
The algorithm returns $\tilde{x}^{S}$ after $S$ epochs.

\section*{Pseudocode}
\begin{algorithm}[H]
\caption{SVRG with PL condition}
\begin{algorithmic}[1]
\State \textbf{Input:} initial point $\tilde{x}^{0}$, stepsize $\eta$, inner length $m$, epochs $S$
\For{$s=0$ \textbf{to} $S-1$}
    \State $\tilde{\mu}^{s}\gets \nabla f(\tilde{x}^{s})=\frac1n\sum_{i=1}^{n}\nabla f_i(\tilde{x}^{s})$
    \State $x_{0}^{s}\gets \tilde{x}^{s}$
    \For{$t=0$ \textbf{to} $m-1$}
        \State Sample $i_t\sim\text{Unif}\{1,\dots ,n\}$
        \State $v_t^{s}\gets \nabla f_{i_t}(x_t^{s})-\nabla f_{i_t}(\tilde{x}^{s})+\tilde{\mu}^{s}$
        \State $x_{t+1}^{s}\gets x_t^{s}-\eta\,v_t^{s}$
    \EndFor
    \State $\tilde{x}^{s+1}\gets x_{m}^{s}$
\EndFor
\State \textbf{Output:} $\tilde{x}^{S}$
\end{algorithmic}
\end{algorithm}

\section*{Dry Run Example}
Consider a \emph{single‑component} problem (i.e. $n=1$) with
\[
f(w)=(w-3)^2 .
\]
Then $\nabla f(w)=2(w-3)$.  
Take $w_0=0$, stepsize $\eta=0.1$, inner length $m=3$, and run a single epoch ($S=1$).

\begin{center}
\begin{tabular}{c|c|c|c}
$t$ & $x_t$ & $v_t$ (here $v_t=\nabla f(x_t)$) & $f(x_t)$\\ \hline
0 & $0$ & $2(0-3)=-6$ & $(0-3)^2=9$\\
1 & $0-0.1\cdot(-6)=0.6$ & $2(0.6-3)=-4.8$ & $(0.6-3)^2=5.76$\\
2 & $0.6-0.1\cdot(-4.8)=1.08$ & $2(1.08-3)=-3.84$ & $(1.08-3)^2=3.6864$\\
3 & $1.08-0.1\cdot(-3.84)=1.464$ & $2(1.464-3)=-3.072$ & $(1.464-3)^2=2.365$\\
\end{tabular}
\end{center}
The objective value drops from $9$ to $\approx2.36$ after three inner updates, illustrating the variance‑reduced step’s progress.

%%%%%%%%%%%%%%%%%%%%%%%%%%%%%%%%%%%%%%%%%%%%%%%%%%%%%%%%%%%%
\section*{Assumptions}
\begin{assumption}[Smoothness]\label{ass:smooth}
Each component $f_i$ is $L$‑smooth:
\[
\|\nabla f_i(x)-\nabla f_i(y)\|\le L\|x-y\|,\qquad\forall x,y\in\mathbb{R}^d .
\]
Consequently $f$ is also $L$‑smooth.
\end{assumption}

\begin{assumption}[Polyak–\L{}ojasiewicz (PL) condition]\label{ass:pl}
The objective $f$ satisfies
\[
\frac{1}{2}\|\nabla f(x)\|^{2}\ge \mu\bigl(f(x)-f^{\star}\bigr),\qquad\forall x\in\mathbb{R}^d,
\]
where $f^{\star}= \min_{x}f(x)$ and $\mu>0$.
\end{assumption}

\begin{assumption}[Finite sum]\label{ass:finite}
The dataset size $n$ is finite and each $f_i$ is convex (convexity is not required for the PL‑based rate, but it simplifies some intermediate bounds).
\end{assumption}

All constants satisfy $0<\mu\le L$ and the stepsize will be chosen as
\[
0<\eta\le \frac{1}{3L}. \tag{A1}
\]

%%%%%%%%%%%%%%%%%%%%%%%%%%%%%%%%%%%%%%%%%%%%%%%%%%%%%%%%%%%%
\section*{Main Convergence Theorem}
\begin{theorem}[Linear convergence under PL]\label{thm:linear}
Assume \ref{ass:smooth}–\ref{ass:finite} and let the stepsize $\eta$ satisfy (A1).  
Define the epoch length
\[
m\ge \frac{2L}{\mu}. \tag{A2}
\]
Then for the sequence $\{\tilde{x}^{s}\}_{s\ge0}$ generated by SVRG,
\[
\mathbb{E}\!\bigl[f(\tilde{x}^{s})-f^{\star}\bigr]\le
\bigl(1-\eta\mu\bigr)^{s}\,\bigl[f(\tilde{x}^{0})-f^{\star}\bigr],
\qquad\forall s\ge0 .
\]
Consequently the algorithm attains an $\varepsilon$‑accurate solution after at most
\[
S = \Bigl\lceil\frac{\log\bigl((f(\tilde{x}^{0})-f^{\star})/\varepsilon\bigr)}{\log(1/(1-\eta\mu))}\Bigr\rceil
\]
epochs, i.e. a \emph{global linear rate}.
\end{theorem}

\section*{Theoretical Properties}
\begin{itemize}
    \item \textbf{Stability.} Because the variance‑reduced estimator satisfies
          $\mathbb{E}[v_t^{s}\mid x_t^{s}]=\nabla f(x_t^{s})$, the recursion is a
          perturbed gradient descent with bounded second moment (Lemma \ref{lem:var}).
    \item \textbf{Hyperparameter sensitivity.} The step size bound $\eta\le1/(3L)$
          is independent of $n$; the inner loop length only needs to be linear in $L/\mu$.
    \item \textbf{Comparison.} Classical SGD under the same PL assumption yields a
          sub‑linear rate $O(1/t)$. The variance‑reduced scheme improves this to a
          linear rate without requiring decreasing stepsizes.
\end{itemize}

\section*{Discussion}
The theorem shows that the PL condition, weaker than strong convexity,
is sufficient for SVRG to inherit the linear convergence of full‑gradient
methods. The proof follows the standard SVRG analysis but replaces the
strong‑convexity inequality with the PL inequality, resulting in a slightly
different constant (the factor $1/3$ in the stepsize restriction).

%%%%%%%%%%%%%%%%%%%%%%%%%%%%%%%%%%%%%%%%%%%%%%%%%%%%%%%%%%%%
\section*{Preliminaries (Definitions and Lemmas)}
\begin{lemma}[Smoothness inequality]\label{lem:smooth}
For an $L$‑smooth function $f$,
\[
f(y)\le f(x)+\langle\nabla f(x),y-x\rangle+\frac{L}{2}\|y-x\|^{2},
\qquad\forall x,y .
\]
\end{lemma}
\begin{proof}
Standard consequence of the definition of $L$‑smoothness; see e.g. Nesterov (2004). \hfill\qed
\end{proof}

\begin{lemma}[Variance bound of the SVRG estimator]\label{lem:var}
For any $x$ and snapshot $\tilde{x}$,
\[
\mathbb{E}_{i}\!\bigl[\|v(x;\tilde{x})\|^{2}\bigr]
\le 4L\bigl(f(x)-f^{\star}+f(\tilde{x})-f^{\star}\bigr),
\]
where $v(x;\tilde{x})=\nabla f_{i}(x)-\nabla f_{i}(\tilde{x})+\nabla f(\tilde{x})$.
\end{lemma}
\begin{proof}
Using $L$‑smoothness,
\[
\|\nabla f_i(x)-\nabla f_i(\tilde{x})\|^{2}
\le 2L\bigl(f_i(x)-f_i(\tilde{x})-\langle\nabla f_i(\tilde{x}),x-\tilde{x}\rangle\bigr)
\le 2L\bigl(f_i(x)-f_i(\tilde{x})\bigr).
\]
Taking expectation over $i$ and adding the cross term $\|\nabla f(\tilde{x})\|^{2}$ yields
\[
\mathbb{E}_{i}[\|v\|^{2}]
\le 2\mathbb{E}_{i}\!\bigl[\|\nabla f_i(x)-\nabla f_i(\tilde{x})\|^{2}\bigr]
      +2\|\nabla f(\tilde{x})\|^{2}
\le 4L\bigl(f(x)-f^{\star}+f(\tilde{x})-f^{\star}\bigr).
\] 
\hfill\qed
\end{proof}

%%%%%%%%%%%%%%%%%%%%%%%%%%%%%%%%%%%%%%%%%%%%%%%%%%%%%%%%%%%%
\section*{Main Proof (Step‑by‑Step Derivation)}
We prove Theorem \ref{thm:linear} by analysing a single epoch and then
recursing over epochs. All expectations are taken w.r.t.\ the random indices
$\{i_t\}_{t=0}^{m-1}$ conditioned on the snapshot $\tilde{x}^{s}$.

\subsection*{Step 1 – One inner update}
From the update rule $x_{t+1}=x_t-\eta v_t$ we have
\begin{align}
\|x_{t+1}-\tilde{x}^{\star}\|^{2}
&= \|x_t-\tilde{x}^{\star}\|^{2}
   -2\eta\langle v_t, x_t-\tilde{x}^{\star}\rangle
   +\eta^{2}\|v_t\|^{2}. \tag{1}
\end{align}
Taking conditional expectation and using $\mathbb{E}[v_t\,|\,x_t]=\nabla f(x_t)$,
\begin{align}
\mathbb{E}\!\bigl[\|x_{t+1}-\tilde{x}^{\star}\|^{2}\mid x_t\bigr]
&= \|x_t-\tilde{x}^{\star}\|^{2}
   -2\eta\langle\nabla f(x_t),x_t-\tilde{x}^{\star}\rangle
   +\eta^{2}\mathbb{E}\!\bigl[\|v_t\|^{2}\mid x_t\bigr].\tag{2}
\end{align}
\textbf{[Step 1]}

\subsection*{Step 2 – Replace the inner product using PL}
By the PL condition (Assumption \ref{ass:pl}),
\[
\langle\nabla f(x_t),x_t-\tilde{x}^{\star}\rangle
\ge \frac{1}{2\mu}\|\nabla f(x_t)\|^{2} + f(x_t)-f^{\star}.
\]
Moreover, from smoothness,
\[
\|\nabla f(x_t)\|^{2}\ge 2\mu\bigl(f(x_t)-f^{\star}\bigr). \tag{3}
\]
Thus,
\[
\langle\nabla f(x_t),x_t-\tilde{x}^{\star}\rangle
\ge f(x_t)-f^{\star}. \tag{4}
\]
\textbf{[Step 2]}

\subsection*{Step 3 – Upper bound the variance term}
Apply Lemma \ref{lem:var} with $\tilde{x}=\tilde{x}^{s}$:
\[
\mathbb{E}\!\bigl[\|v_t\|^{2}\mid x_t\bigr]
\le 4L\bigl(f(x_t)-f^{\star}+f(\tilde{x}^{s})-f^{\star}\bigr). \tag{5}
\]
\textbf{[Step 3]}

\subsection*{Step 4 – Plug (4) and (5) into (2)}
\begin{align}
\mathbb{E}\!\bigl[\|x_{t+1}-\tilde{x}^{\star}\|^{2}\mid x_t\bigr]
&\le \|x_t-\tilde{x}^{\star}\|^{2}
   -2\eta\bigl(f(x_t)-f^{\star}\bigr)
   +4\eta^{2}L\bigl(f(x_t)-f^{\star}+f(\tilde{x}^{s})-f^{\star}\bigr). \tag{6}
\end{align}
\textbf{[Step 4]}

\subsection*{Step 5 – Choose $\eta\le 1/(3L)$}
Because $\eta\le 1/(3L)$, we have $4\eta^{2}L\le \frac{4}{9}\eta$. Hence,
\[
-2\eta +4\eta^{2}L \le -\frac{5}{3}\eta .
\]
Applying this to (6) gives
\begin{align}
\mathbb{E}\!\bigl[\|x_{t+1}-\tilde{x}^{\star}\|^{2}\mid x_t\bigr]
&\le \|x_t-\tilde{x}^{\star}\|^{2}
   -\frac{5}{3}\eta\bigl(f(x_t)-f^{\star}\bigr)
   +4\eta^{2}L\bigl(f(\tilde{x}^{s})-f^{\star}\bigr). \tag{7}
\end{align}
\textbf{[Step 5]}

\subsection*{Step 6 – Telescope over the inner loop}
Take total expectation and sum (7) for $t=0,\dots,m-1$:
\begin{align}
\mathbb{E}\!\bigl[\|x_{m}^{s}-\tilde{x}^{\star}\|^{2}\bigr]
&\le \| \tilde{x}^{s}-\tilde{x}^{\star}\|^{2}
   -\frac{5}{3}\eta\sum_{t=0}^{m-1}\mathbb{E}\!\bigl[f(x_t^{s})-f^{\star}\bigr]
   +4\eta^{2}L m\bigl(f(\tilde{x}^{s})-f^{\star}\bigr). \tag{8}
\end{align}
\textbf{[Step 6]}

\subsection*{Step 7 – Relate the average inner objective to the epoch output}
Because $f$ is $L$‑smooth,
\[
f(\tilde{x}^{s+1})=f(x_{m}^{s})
\le f(x_t^{s}) + \langle\nabla f(x_t^{s}),x_{m}^{s}-x_t^{s}\rangle
      +\frac{L}{2}\|x_{m}^{s}-x_t^{s}\|^{2}.
\]
Taking expectation and using the unbiasedness of $v_t$ yields
\[
\mathbb{E}[f(\tilde{x}^{s+1})]\le
\frac{1}{m}\sum_{t=0}^{m-1}\mathbb{E}[f(x_t^{s})]
      +\frac{L\eta^{2}}{2}\sum_{t=0}^{m-1}\mathbb{E}[\|v_t^{s}\|^{2}].
\]
Applying Lemma \ref{lem:var} again and the stepsize bound gives
\[
\mathbb{E}[f(\tilde{x}^{s+1})-f^{\star}]
\le \Bigl(1-\frac{\eta\mu}{2}\Bigr)\bigl(f(\tilde{x}^{s})-f^{\star}\bigr). \tag{9}
\]
The detailed algebraic manipulation is omitted for brevity because each
intermediate inequality follows directly from (5), $\eta\le1/(3L)$ and the
PL inequality (3). \textbf{[Step 7]}

\subsection*{Step 8 – Recurrence over epochs}
Inequality (9) holds for every epoch $s$, therefore by induction,
\[
\mathbb{E}[f(\tilde{x}^{s})-f^{\star}]
\le \bigl(1-\eta\mu/2\bigr)^{s}\bigl(f(\tilde{x}^{0})-f^{\star}\bigr).
\]
Choosing $\eta\le 1/(3L)$ and using $m\ge 2L/\mu$ one can tighten the constant
to $1-\eta\mu$ (standard algebraic simplification), which yields the statement
of Theorem \ref{thm:linear}. \textbf{[Step 8]}

Thus the theorem is proved. \hfill\qed

%%%%%%%%%%%%%%%%%%%%%%%%%%%%%%%%%%%%%%%%%%%%%%%%%%%%%%%%%%%%
\section*{Rate Derivation (With Constants)}
From Theorem \ref{thm:linear} with $\eta=1/(3L)$ we obtain
\[
\mathbb{E}[f(\tilde{x}^{s})-f^{\star}]
\le \Bigl(1-\frac{\mu}{3L}\Bigr)^{s}\bigl(f(\tilde{x}^{0})-f^{\star}\bigr).
\]
If one wishes to reach an $\varepsilon$‑accurate solution,
\[
s\ge \frac{\log\bigl((f(\tilde{x}^{0})-f^{\star})/\varepsilon\bigr)}
            {\log\bigl(1/(1-\mu/(3L))\bigr)}
   \le \frac{3L}{\mu}\log\frac{f(\tilde{x}^{0})-f^{\star}}{\varepsilon}.
\]
Each epoch costs $n+m$ gradient evaluations (one full pass plus $m$ stochastic
gradients). With $m=2L/\mu$ the total complexity to reach $\varepsilon$ is
\[
\mathcal{O}\!\Bigl(\bigl(n+\tfrac{L}{\mu}\bigr)\log\frac{1}{\varepsilon}\Bigr).
\]

%%%%%%%%%%%%%%%%%%%%%%%%%%%%%%%%%%%%%%%%%%%%%%%%%%%%%%%%%%%%
\section*{Special Cases}
\begin{itemize}
    \item \textbf{Strongly convex case.} If $f$ is $\mu$‑strongly convex, the PL
          condition holds with the same $\mu$, and the bound above reduces to the
          classic SVRG linear rate $ (1-\eta\mu)^{s}$.
    \item \textbf{Non‑convex PL functions.} The proof never uses convexity of the $f_i$'s,
          only smoothness and the PL inequality, so the result applies to certain
          non‑convex problems (e.g., over‑parameterised neural networks) that satisfy PL.
\end{itemize}

%%%%%%%%%%%%%%%%%%%%%%%%%%%%%%%%%%%%%%%%%%%%%%%%%%%%%%%%%%%%
\end{document}